% ============================================================================
% Chapter 4 — Analysis and Results
% Figures 1-7 are numbered in this section, in order, to match the companion
% code (dlvt/figures.py) and the test-suite docstrings.
% ============================================================================
\section{Analysis and Results}
\label{sec:results}

This section reports the phase-plane analysis at the baseline calibration,
the comparative statics---including the corrected reading of the
$\beta$-sweep---and the falsifiable proposition set. Every number quoted
here is pinned by an automated test in the companion repository.

\subsection{Baseline dynamics}
\label{sec:results-baseline}

Figure~\ref{fig:temporal} traces the baseline trajectory from
$(V_0,C_0)=(8.0,5.0)$ and shows the model's causal loop unfolding in time.
Three features should be visible. First, an early benign phase: complexity
is low, the depletion ratio $\Gamma\ll1$, and vitality holds near the
ceiling while capital compounds. Second, the turn: as capital accumulates,
complexity---and with it the drain term of \eqref{eq:dVdt}---grows
quadratically ($\gamma=2$), and vitality is pulled down through the
strategic threshold even as capital continues to rise. Third, convergence:
the system spirals into the unique interior equilibrium
$(V^{*},C^{*})=(4.7025,\,32.0337)$ with eigenvalues $-0.205\pm0.331i$
(rotation period $\approx19$ time units, amplitude e-folding time
$\approx4.9$; Appendix~\ref{app:jacobian}). The transient overshoot is
diagnostic rather than incidental: from $(10,40)$ capital peaks near
$C\approx46.4$, above the carrying capacity $\Cmax=44.99$, before
complexity friction pulls it back---visual confirmation that $\Cmax$
bounds sustainable flux, not reachable states (Remark~\ref{rem:ctrap}).

\begin{figure}[t]
\centering
\includegraphics[width=0.95\linewidth]{fig1_temporal_evolution.pdf}
\caption{Temporal evolution of $V(t)$, $C(t)$, $O(t)$, and $\Gamma(t)$ at
the baseline calibration (Table~\ref{tab:parameters}), $(V_0,C_0)=(8,5)$.}
\label{fig:temporal}
\end{figure}

Figure~\ref{fig:scenarios} places the three outcome regimes of
Definition~\ref{def:lowvitality} side by side: three qualitatively
different vitality trajectories from identical initial conditions. In the
sustainable panel (low coupling and low cost, $\delta=0.008$,
$\beta=0.15$) vitality settles above $\Vstrat$; in the baseline panel it
settles below; and in the collapse-prone panel (parameters with
$\Cmax=0$) even $C=0$ implies $\Gamma>1$, and vitality is driven into the
depletion band (Proposition~\ref{prop:band}(ii)). The figure establishes
that the regime label is a property of parameters, not of starting
points---the claim that Propositions P1 and P2 below make testable.

\begin{figure}[t]
\centering
\includegraphics[width=0.95\linewidth]{fig2_three_scenarios.pdf}
\caption{Three outcome regimes: sustainable, low-vitality, collapse-prone.}
\label{fig:scenarios}
\end{figure}

Figure~\ref{fig:phase} compresses the same dynamics into the phase plane.
Two curves organize the portrait: the strictly decreasing vitality
nullcline (Proposition~\ref{prop:nullcline}) and the increasing capital
nullcline \eqref{eq:cnullcline}. Their unique interior crossing
(Theorem~\ref{thm:uniqueness}) is the attractor, and the flow field shows
every plotted trajectory spiraling into it---no second equilibrium, no
closed orbit, no basin boundary anywhere in the frame. This is the
single-attractor geometry that separates DLVT from bistable capability
traps.

\begin{figure}[t]
\centering
\includegraphics[width=0.8\linewidth]{fig3_phase_portrait.pdf}
\caption{Phase portrait with nullclines and the unique interior equilibrium
$(V^{*},C^{*})\approx(4.70,\,32.03)$.}
\label{fig:phase}
\end{figure}

\subsection{Comparative statics and the corrected $\beta$-sweep}
\label{sec:results-statics}

Figure~\ref{fig:bifurcation} sweeps $\beta$ and $R$, and the two panels
should look strikingly different. The $R$-panel behaves as intuition
suggests: equilibrium vitality rises with recovery capacity. The
$\beta$-panel carries the substantive result: $V^{*}(\beta)$ is
\emph{flat}, while $C^{*}(\beta)$ falls along the hyperbola
$C^{*}(\beta)=(O^{*}-O_0)/\beta=8.008/\beta$ (at $\varepsilon=0.1$;
$7.933$ in the $\varepsilon\to0$ limit). By Corollary~\ref{cor:scope},
$(V^{*},O^{*})$ solve a $\beta$-free system, so all $\beta$-dependence is
absorbed into capital. Two consequences follow.

First, there is \emph{no bifurcation in $\beta$} within the power-law
family. An earlier draft reported a critical coupling
$\beta_{\mathrm{crit}}\approx0.1015$ at which the equilibrium appeared to
cross the strategic threshold; that value was a scan-window artifact---the
equilibrium branch $C^{*}(\beta)=8.008/\beta$ exits a legacy search window
$C\le80$ precisely at $\beta\approx0.100$---and it is retracted. The full
forensic account, together with the honest interval-reporting protocol now
implemented in the companion code, is Appendix~\ref{app:bifurcation}.

Second, the flatness of $V^{*}(\beta)$ is the formal content of the
\emph{intervention asymmetry}: reducing the scope coupling reallocates the
equilibrium between capital and complexity but leaves equilibrium vitality
untouched, whereas recovery-side parameters move $V^{*}$ directly.

\begin{figure}[t]
\centering
\includegraphics[width=0.95\linewidth]{fig4_bifurcation.pdf}
\caption{Equilibrium loci under parameter sweeps: $C^{*}$ and $V^{*}$
against $\beta$ and $R$. Note the $\beta$-invariance of $V^{*}$
(Corollary~\ref{cor:scope}) and the conserved product $\beta C^{*}$.}
\label{fig:bifurcation}
\end{figure}

Figure~\ref{fig:impact} makes the deployment wedge visible: it plots
realized impact under DLVT against the counterfactual in which the same
capital is deployed without an energy gate---the implicit assumption of
classical human-capital reasoning \citep{becker1993human}. The two curves
coincide early, while vitality is high and the gate is effectively open,
then separate as complexity-driven drain lowers vitality; near equilibrium
the ungated counterfactual continues to credit the leader with the full
yield of a large capital stock while realized impact plateaus well below
it. The gap between the curves is the model's central object: capital
owned but not usable.

\begin{figure}[t]
\centering
\includegraphics[width=0.9\linewidth]{fig5_impact_comparison.pdf}
\caption{Realized impact under DLVT versus the energy-ungated
human-capital counterfactual.}
\label{fig:impact}
\end{figure}

\subsection{Carrying capacity and the regime map}
\label{sec:results-regimes}

Figure~\ref{fig:heatmap} maps the carrying capacity $\Cmax(\beta,R)$ of
Proposition~\ref{prop:carrying}. The level sets should be read as
comparative statics made visible: $\Cmax$ falls with the coupling $\beta$
and the energetic cost $\delta$, and rises with recovery capacity $R$, per
\eqref{eq:cmax-statics}. We reiterate the corrected interpretation: $\Cmax$
is the flux threshold at which $\Gamma=1$ at full vitality, not a fold
point---$\det J>0$ everywhere (Lemma~\ref{lem:trace}), so no equilibria
are created or destroyed there, and trajectories may transiently exceed
it.

Figure~\ref{fig:regimemap} partitions the $(\beta,\delta)$ plane into the
three regimes of Definition~\ref{def:lowvitality}; the reader should note
how close the baseline point sits to the sustainable/low-vitality
boundary. That proximity is quantified by the regime governor: at fixed
drain-side parameters the label turns on $\mu/\alpha$, and the baseline
$\mu/\alpha=2.0$ sits only $8\%$ below the flip value $2.163$ at which
$V^{*}=\Vstrat$. Across a $\pm2\times$ log-uniform hypercube around
baseline, $P(\text{low-vitality}\mid\text{stable
equilibrium})\approx0.49$: the illustrative calibration lies near the
center of its own regime boundary, which is the honest way to read every
baseline classification in this paper.

\begin{figure}[t]
\centering
\includegraphics[width=0.85\linewidth]{fig6_sensitivity_heatmap.pdf}
\caption{Carrying capacity $\Cmax(\beta,R)$; comparative statics per
Eq.~\eqref{eq:cmax-statics}.}
\label{fig:heatmap}
\end{figure}

\begin{figure}[t]
\centering
\includegraphics[width=0.85\linewidth]{fig7_regime_map.pdf}
\caption{Regime classification in the $(\beta,\delta)$ plane
(Definition~\ref{def:lowvitality}).}
\label{fig:regimemap}
\end{figure}

\subsection{Falsifiable propositions}
\label{sec:propositions}

We close the analysis with the proposition set the theory stakes its claim
on. These are \emph{empirical} propositions, labeled P1--P6 to distinguish
them from the formal propositions of \S\ref{sec:formulation}, and each is
stated with the observation that would falsify it. Before stating them, we
declare where they apply. The boxed scope conditions below are part of the
theory's content, not fine print: they specify the leaders, roles, and
observation regimes for which P1--P6 are claims at all, and they follow
directly from the assumptions and robustness analysis of
\S\ref{sec:formulation}.

\begin{center}
\fbox{\begin{minipage}{0.92\linewidth}
\textbf{Scope conditions for P1--P6.} The propositions apply to individual
leaders in roles where (i) complexity load is predominantly generated by the
leader's own accumulated scope (Assumption~\ref{ass:scope}); (ii) the
load--drain relationship is convex ($\gamma>1$, \S\ref{sec:robustness});
(iii) the drain kernel is well approximated by the separable power-law
family (P3--P4 delimit this boundary); and (iv) observation windows are long
enough for equilibrium behavior ($\gtrsim$ several e-folding times;
$\approx5$ time units at baseline). Outside these conditions the theory is
silent, not merely weakened.
\end{minipage}}
\end{center}

\begin{description}
\item[P1 (Single attractor).] We propose that individual vitality--capital
trajectories converge to a single interior steady state, with
damped-spiral transients. \emph{Falsifier:} persistent oscillation, or
evidence of multistability---distinct long-run outcomes for similar
leaders in similar roles that are attributable to initial conditions
rather than parameters.

\item[P2 (Regime governor).] We propose that the low-vitality regime
occurs if and only if the depreciation--accumulation ratio $\mu/\alpha$ is
small relative to the threshold implied by \eqref{eq:vstar}.
\emph{Falsifier:} regime membership uncorrelated with the
depreciation/accumulation ratio across leaders.

\item[P3 (Scope absorption).] We propose that reducing the scope coupling
$\beta$ alone does not durably raise equilibrium vitality: capital
re-expands until the product $\beta C^{\eta}$ is restored.
\emph{Falsifier:} a durable vitality gain from coupling reduction alone,
with $\beta C^{*}$ \emph{not} conserved.

\item[P4 (Boundary of P3).] We propose that the magnitude of
$\dd V^{*}/\dd\beta$ grows with the non-separability of the empirical
drain kernel, so that where drain saturates in complexity (Hill-type),
scope reduction regains leverage. This proposition marks the edge of P3's
validity rather than extending it. \emph{Falsifier:} scope-reduction
effects unrelated to the curvature of the estimated drain kernel.

\item[P5 (Carrying-capacity statics).] We propose that the maximum
sustainable capital $\Cmax$ falls with the coupling $\beta$ and the
energetic cost $\delta$, and rises with recovery capacity $R$, per
\eqref{eq:cmax-statics}. \emph{Falsifier:} observed sustainable-capital
ceilings that do not move with recovery capacity, or that rise with
coupling.

\item[P6 (Intervention ordering).] We propose that recovery-side
interventions---raising $R$, lowering $\mu$, raising $\alpha$---raise
equilibrium vitality, while scope redesign (lowering $\beta$) does not.
\emph{Falsifier:} a well-identified field comparison in which scope
redesign outperforms recovery support on long-run vitality. This is the
theory's killer test (\S\ref{sec:roadmap}).
\end{description}
